% =====================================================================
\section{高扭度：\texorpdfstring{$1/Q^2$}{1/Q2} 处的广义因子化及其在更高幂次的失效}\label{sec:dy-twist}
% =====================================================================

\cref{thm:dy-factorization}的余项 $\order{(\Lambda/Q)^p}$ 中 $p$ 不由证明确定。对 DIS，幂次修正由 OPE 的高扭度算符给出（\cref{prop:dis-twist4}，\tagR，$R_{\mathrm{OPE}}$）；DY 不是 OPE 可及的量，幂次修正只能在动量空间中逐区域分析。已知的情形是：相对领头幂次的第一阶修正（$1/Q^2$，扭度 4）仍有广义的因子化形式，第二阶（$m^4/s^2$，扭度 6）则没有。

\begin{theorem}[$1/Q^2$ 处的广义因子化 \tagA]\label{thm:dy-higher-twist}
DY 截面的领头幂次修正是相对 $1/Q^2$ 的扭度 4 项，具有广义因子化形式：它们是一个可计算的短程硬部分、一个强子的扭度 2 矩阵元（部分子分布）与另一个强子的扭度 4 多部分子关联矩阵元的卷积，
\[
\frac{\dd\sigma}{\dd Q^2}=\frac{\dd\sigma^{(\rm tw\,2)}}{\dd Q^2}+\frac{1}{Q^2}\sum_{i,j}\Big[f_{i/A}\otimes T^{(4)}_{j/B}\otimes H^{(4)}_{ij}+(A\leftrightarrow B)\Big]+\order{\frac{1}{Q^4}},
\]
其中卷积是示意的写法（$T^{(4)}$ 依赖于多个动量分数，卷积作用于全部动量分数），$T^{(4)}_{j/B}$ 是含四个场（例如一对夸克场与一对场强，或两对夸克场）的光锥算符在强子 $B$ 中的矩阵元\cite{QiuSterman1991a,QiuSterman1991b}。同样的结构用于原子核中的双散射\cite{Kastella:1989ux}。\\
非微扰地位：\cref{asm:intro-bridge}（MA1、MA2）。
% errtype: E-S
\uses{thm:dy-factorization,asm:intro-bridge}
\end{theorem}

\begin{proofsketch}
(1) 在次领头幂次上，额外的一对共线部分子（夸克或物理极化的胶子）从一个强子的共线子图连到硬部分，按幂次计数每多一对代价为 $\Lambda^2/Q^2$；纵向胶子仍由 Ward 恒等式收成 Wilson 线，使多部分子关联函数规范不变。(2) 末态相互作用的幺正性相消与软胶子的处理在这一幂次上重复领头幂次的论证；在标量夸克模型中，单圈软胶子效应在 $m^2/s$ 处确实因子化\cite{Basu:1984ba}。(3) 由此得到式中的三重卷积，硬部分 $H^{(4)}$ 可逐阶计算。Qiu 与 Sterman 在 $\as$ 的领头阶算出了这些修正，并论证了这类 $1/Q^2$ 项可以由广义因子化定理处理\cite{QiuSterman1991a,QiuSterman1991b}（据其摘要）。（各幂次计数步骤只给出 E-S 型陈述。）
\end{proofsketch}

\paragraph{地位与缺口}\tagA。扭度 4 的全阶论证见\cite{QiuSterman1991b}；\cref{app:appO}的 Glauber 分析只在领头幂次进行，次领头幂次上的单侧改写与形变没有写成条件定理，这是本定理比\cref{thm:dy-factorization}更弱的地方。CSS 的综述指出，微扰论在 $m^2/s$ 处遵守因子化，这使得因子化定理在这一幂次上对强子—强子散射可能成立\cite[第~10.2~节]{CollinsSoperSterman1989}。还要注意，领头扭度微扰级数的红外重整化子不确定性与 $1/Q^2$ 修正同阶（\cref{asm:intro-bridge}(Br3)；对 DY 见\cite{BenekeBraun1995b}），因此扭度 4 矩阵元只有在规定了微扰级数的求和方式之后才有确定的意义。

\begin{proposition}[更高幂次的因子化失效 \tagV]\label{prop:dy-twist6-failure}
在 $m^4/s^2$（相对领头幂次的第二阶幂次修正，扭度 6）上，Drell--Yan 截面不具有广义因子化形式：在微扰论中，这一幂次上存在既不在实、虚修正之间相消，也不能吸收进强子矩阵元的红外敏感贡献；而在 $m^2/s$ 上因子化仍成立（\cref{thm:dy-higher-twist}）。
\uses{thm:dy-higher-twist}
\end{proposition}

\paragraph{地位与缺口}\tagV。证据有三类。(i) 模型中的显式微扰反例：在 Drell--Yan 过程的标量夸克 QCD 模型中，单圈软胶子效应在 $m^2/s$ 处因子化，在 $m^4/s^2$ 处不因子化，而且非因子化项不被形状因子压低\cite{Basu:1984ba}（据其摘要）。(ii) QCD 中的显式计算：CSS 的综述总结道，在高扭度上强子—强子散射的红外发散超出单圈即不相消，在 Drell--Yan 中两圈的非相消发散出现在 $m^4/s^2$ 量级\cite[第~10.2~节]{CollinsSoperSterman1989}（另见\cref{rem:dy-dft-premises}）。(iii) Bodwin、Braaten 与 Lepage 的总结：在 $1/Q^2$ 的次领头阶之外，“factorization is known to fail”\cite[第~VI~节]{BodwinBraatenLepage1995}。这些都是微扰论的陈述。CSS 同时强调，如何解释这种不相消“is not quite clear”，包括非微扰效应在内的实际情况可能比微扰计算所显示的更好或更坏\cite[第~10.2~节]{CollinsSoperSterman1989}；因此 \tagV 指的是微扰论中广义因子化形式的失效，它对强子数据的数值影响并未确定。
